\documentclass[a4paper, serif, 10pt]{moderncv}

%% moderncv
% style and color
\moderncvstyle{banking}
\moderncvcolor{black}

%% page layout/margins
\usepackage[top=2.5cm, bottom=2.5cm, left=1.5cm, right=1.5cm]{geometry}

\nopagenumbers{}

%% Babel config for French language
% ref:
% - https://latex3.github.io/babel/guides/locale-french.html
\usepackage[french]{babel}

%% fontspec
\usepackage{fontspec}

\IfFontExistsTF{Linux Libertine}{
  \setmainfont[Ligatures={TeX, Common}, Numbers=OldStyle]{Linux Libertine}
}{}
\IfFontExistsTF{Linux Libertine O}{
  \setmainfont[Ligatures={TeX, Common}, Numbers=OldStyle]{Linux Libertine O}
}{}

%% CTeX
% CJK support, used to show CJK characters in the resume
%
% - fontset=none: disable builtin fontset but instead set the CJK font manually
% - heading=false: disable ctex heading
% - punct=kaiming: use kaiming punctuations styles for CJK
% - scheme=plain: use plain scheme, do not override \normalsize font size
% - space=auto: space settings for CJK characters
%
% ref:
% - http://ctan.mirrorcatalogs.com/language/chinese/ctex/ctex.pdf
\usepackage[UTF8, heading=false, punct=kaiming, scheme=plain, space=auto]{ctex}

\IfFontExistsTF{Noto Serif CJK SC}{
  \setCJKmainfont{Noto Serif CJK SC}
}{}
\IfFontExistsTF{Noto Sans CJK SC}{
  \setCJKsansfont{Noto Sans CJK SC}
}{}

%% line spacing
\usepackage{setspace}
\setstretch{1.125}

%% URL styling - use normal text instead of monospace
\urlstyle{same}

% Auto-underline all links
% Defer until \begin{document} so hyperref is already loaded
\AtBeginDocument{%
  \let\oldhref\href
  \renewcommand{\href}[2]{\oldhref{#1}{\underline{#2}}}%
}

%% Patch \httplink and \httpslink to handle full URLs
% The original moderncv macros always prepend http:// or https:// to URLs.
% This causes issues when the URL already has a protocol (e.g., https://example.com)
% resulting in malformed URLs like https://https://example.com.
% This patch checks if the URL already contains "://" and uses it directly if so.
% Note: We use \str_set:Nx (x-expansion) to expand macros like \@homepage first.
\ExplSyntaxOn
\str_new:N \g_yamlresume_url_str

\RenewDocumentCommand{\httpslink}{O{}m}{
  \str_set:Nx \g_yamlresume_url_str {#2}
  \str_if_in:NnTF \g_yamlresume_url_str {://}
    { \url{#2} }
    { \url{https://#2} }
}

\RenewDocumentCommand{\httplink}{O{}m}{
  \str_set:Nx \g_yamlresume_url_str {#2}
  \str_if_in:NnTF \g_yamlresume_url_str {://}
    { \url{#2} }
    { \url{http://#2} }
}
\ExplSyntaxOff

\name{Camille Dubois}{}
\title{Ingénieur DevOps / SRE}
\phone[mobile]{+33 6 12 34 56 78}
\email{camille.dubois@protonmail.com}
\homepage{https://camilledubois.dev}

\address{12 rue de la Roquette, Paris, Île-de-France, France, 75011}{}{}

\extrainfo{{\small \faGithub}\ \href{https://github.com/camilledubois}{@camilledubois} {} {} {} • {} {} {} 
{\small \faLinkedin}\ \href{https://www.linkedin.com/in/camille-dubois-devops}{@camille-dubois-devops}}

\begin{document}

\maketitle

\section{Informations de base}

\cvline{}{Ingénieur DevOps avec plus de 5 ans d'expérience dans la conception, l'automatisation et l'exploitation de plateformes cloud natives. Passionné par la fiabilité, l'observabilité et l'amélioration continue.
%
\begin{itemize}
\item Expertise en Kubernetes, Terraform et CI/CD sur AWS et Azure
\item Réduction des coûts cloud de 30 \% grâce à l'optimisation des ressources et au rightsizing
\item Mise en place de pipelines de déploiement sans interruption pour plus de 200 microservices
\item Solide culture SRE : SLO, gestion des incidents, post-mortems et ingénierie du chaos
\end{itemize}}
\section{Formation}

\cventry{sept. 2019–juil. 2021}
        {Master, Informatique - Systèmes distribués et cloud, Score : Mention Très Bien}
        {Université Paris-Saclay}
        {\url{https://www.universite-paris-saclay.fr}}
        {}
        {\begin{itemize}
\item Spécialisation en systèmes distribués, cloud computing et automatisation
\item Projet de fin d'études : orchestration Kubernetes multi-cloud avec Terraform
\item Major de promotion en architecture cloud et sécurité
\end{itemize}
\textbf{Cours} : Architectures cloud et virtualisation, Systèmes distribués et tolérance aux pannes, Sécurité des infrastructures, Automatisation et Infrastructure as Code, Observabilité et performance}

\cventry{sept. 2016–juin 2019}
        {Licence, Informatique, Score : Mention Bien}
        {Université de Lille}
        {\url{https://www.univ-lille.fr}}
        {}
        {\begin{itemize}
\item Acquisition des fondamentaux en développement, réseaux et systèmes d'exploitation
\item Premier contact avec les pratiques DevOps lors de projets universitaires
\end{itemize}
\textbf{Cours} : Algorithmique et structures de données, Réseaux et protocoles, Administration système Linux, Programmation Python et Bash}
\section{Expérience professionnelle}

\cventry{janv. 2023–Present}
        {Ingénieur DevOps Senior}
        {CloudNova}
        {\url{https://www.cloudnova.fr}}
        {}
        {\begin{itemize}
\item Conception et maintien d'une plateforme Kubernetes multi-tenant sur AWS (EKS) pour 200+ microservices
\item Automatisation complète du provisionnement avec Terraform et Ansible, réduisant le time-to-market de 40 \%
\item Mise en place de pipelines GitLab CI/CD avec tests automatisés, scans de sécurité et déploiements canary
\item Définition et suivi des SLO/SLI, animation des astreintes et des post-mortems
\item Optimisation des coûts cloud de 30 \% via rightsizing, spot instances et politiques de rétention
\end{itemize}
\textbf{Mots-clés} : Kubernetes, Terraform, AWS, GitLab CI, SRE, Observabilité}

\cventry{sept. 2021–déc. 2022}
        {Ingénieur DevOps}
        {PayFlow}
        {\url{https://www.payflow.fr}}
        {}
        {\begin{itemize}
\item Migration d'une application monolithique vers une architecture de microservices conteneurisés
\item Création de charts Helm et de pipelines Jenkins pour des déploiements reproductibles
\item Mise en place de la supervision avec Prometheus, Grafana et Alertmanager
\item Participation à la rotation d'astreinte et à la réduction du MTTR de 50 \%
\end{itemize}
\textbf{Mots-clés} : Docker, Helm, Jenkins, Prometheus, Grafana}

\cventry{juil. 2019–août 2021}
        {Administrateur Systèmes et Réseaux}
        {Startup Factory}
        {\url{https://www.startupfactory.fr}}
        {}
        {\begin{itemize}
\item Administration de serveurs Linux (Debian/Ubuntu) et gestion des sauvegardes
\item Automatisation des tâches récurrentes avec Bash et Python
\item Mise en place de la supervision Nagios et du référencement des incidents
\item Support technique aux équipes de développement et participation aux revues d'architecture
\end{itemize}
\textbf{Mots-clés} : Linux, Bash, Python, Nagios, Réseaux}
\section{Langues}

\cvline{Français}{Niveau natif ou bilingue \hfill \textbf{Mots-clés} : Langue maternelle}
\cvline{Anglais}{Niveau professionnel complet \hfill \textbf{Mots-clés} : TOEIC 950, Anglais technique}
\cvline{Espagnol}{Niveau de travail limité \hfill \textbf{Mots-clés} : Conversation courante}
\section{Compétences}

\cvline{Cloud \& Infrastructure}{Expert \hfill \textbf{Mots-clés} : AWS, Azure, GCP, Terraform, Ansible, Pulumi}
\cvline{Conteneurs \& Orchestration}{Expert \hfill \textbf{Mots-clés} : Kubernetes, Docker, Helm, Istio, EKS, Kustomize}
\cvline{CI/CD \& Automatisation}{Avancé \hfill \textbf{Mots-clés} : GitLab CI, GitHub Actions, Jenkins, ArgoCD, Python, Go}
\cvline{Observabilité \& SRE}{Avancé \hfill \textbf{Mots-clés} : Prometheus, Grafana, ELK, Datadog, SLO/SLI, Ingénierie du chaos}
\cvline{Sécurité \& Conformité}{Intermédiaire \hfill \textbf{Mots-clés} : Vault, Trivy, OPA, RBAC, SOC2, RGPD}
\section{Récompenses}

\cventry{juin 2023}
        {CloudNova}
        {Prix de l'innovation DevOps}
        {}
        {}
        {Récompense interne pour la conception d'une plateforme d'auto-scaling prédictif ayant permis une réduction de 25 \% des coûts d'infrastructure.}
\section{Certificats}

\cventry{mars 2023}
        {Amazon Web Services}
        {AWS Certified DevOps Engineer – Professional}
        {\url{https://aws.amazon.com/certification/}}
        {}
        {}

\cventry{nov. 2022}
        {Cloud Native Computing Foundation}
        {Certified Kubernetes Administrator (CKA)}
        {\url{https://www.cncf.io/certification/cka/}}
        {}
        {}

\cventry{mai 2022}
        {HashiCorp}
        {HashiCorp Certified: Terraform Associate}
        {\url{https://www.hashicorp.com/certification}}
        {}
        {}
\section{Publications}

\cventry{avr. 2024}
        {Retour d'expérience : migration de 200 microservices vers Kubernetes}
        {Blog technique de CloudNova}
        {\url{https://www.cloudnova.fr/blog/migration-kubernetes}}
        {}
        {\begin{itemize}
\item Présentation des défis rencontrés lors de la migration d'un monolithe vers Kubernetes
\item Bonnes pratiques pour la gestion des secrets, des ressources et des déploiements progressifs
\item Analyse des gains de performance et de fiabilité après migration
\end{itemize}}
\section{Projets}

\cventry{févr. 2023–août 2023}
        {Mise en place d'une stack d'observabilité centralisée pour 200+ microservices}
        {Plateforme d'observabilité interne}
        {\url{https://github.com/camilledubois/observability-platform}}
        {}
        {\begin{itemize}
\item Déploiement de Prometheus, Grafana et Loki sur Kubernetes
\item Création de tableaux de bord SLO et d'alertes basées sur les symptômes
\item Réduction du MTTR de 50 \% et amélioration de la visibilité sur les incidents
\end{itemize}
\textbf{Mots-clés} : Observabilité, Prometheus, Grafana, Loki, SLO}

\cventry{mars 2022–nov. 2022}
        {Provisionnement automatisé d'environnements éphémères sur AWS et Azure}
        {Automatisation multi-cloud}
        {\url{https://github.com/camilledubois/multi-cloud-automation}}
        {}
        {\begin{itemize}
\item Développement de modules Terraform réutilisables et de pipelines GitLab CI
\item Intégration de politiques OPA et de scans de sécurité Trivy
\item Réduction du temps de création d'un environnement de 2 jours à 30 minutes
\end{itemize}
\textbf{Mots-clés} : Terraform, AWS, Azure, OPA, GitLab CI}
\section{Centres d'intérêt}

\cvline{Open Source}{Kubernetes, Terraform, Prometheus}
\cvline{Sport}{Course à pied, Escalade}
\cvline{Lecture}{Essais techniques, Science-fiction}
\section{Bénévolat}

\cventry{sept. 2022–Present}
        {Mentor}
        {DevOps Paris}
        {\url{https://www.meetup.com/devops-paris/}}
        {}
        {\begin{itemize}
\item Accompagnement de développeurs et d'administrateurs dans leur transition vers les pratiques DevOps
\item Animation d'ateliers pratiques sur Kubernetes, Terraform et CI/CD
\item Organisation de rencontres mensuelles et de retours d'expérience
\end{itemize}}
    
\end{document}